% Transcription — fonds Alexandre Grothendieck, shelfmark 43, batch 11 (pages 201-207).
% Established from the facsimile archives/batches/43/batch-11.pdf (cut from a
% copy of 43.pdf fetched through the project's relay,
% https://grothendieck-relay.onrender.com/source/43.pdf; PDF page k =
% archivists' page 200+k, no cover sheet), per the skill
% transcribe-grothendieck. The folder is 207 pages in eleven batches; this is
% the last, seven pages. Batches 9 and 10 were transcribed in parallel; both
% files existed when this one was finished and were read: their notation is
% followed for the Dieudonné run of page 201 (A the Dieudonné ring, F and V,
% categories C_1, ...), and folders 59 and 9 for the rest
% (\underline{\mathrm{Pic}}, NS, \mathbb{G}_m).
%
% Dating: the folder has its own inventory date, carried verbatim in
% \dating{} with no group sentence, as in batches 9 and 10. No page of this
% batch carries a date.
%
% Pages skipped: none. Page 202 is a cover (an orange sheet inscribed in his
% hand at the top right, « Accouplements de [one sign] Tate -- Lang ») and
% takes no \page marker; its words become the section heading, with a \note.
%
% Pages summarised: none.
%
% Copies. The folder is « notes et copies de notes manuscrites ». On these
% seven pages nothing can be identified with confidence as a photocopy:
% page 201 is in blue ink on the back of a varitypist's printed letterhead
% (upside down, not transcribed); pages 203-204 are black ink on grey paper,
% page 203 headed in orange felt pen « Accouplements de Tate » — coloured ink,
% so at least that heading is not a copy; pages 205-207 are ink on yellowed
% paper. All are transcribed as originals. Whoever checks against the leaves
% should settle 203-204 first.
%
% His pagination: none visible on pages 201-207. Page 201 opens mid-sentence
% (« m/m^2 dual de P(U) primitifs, ... ») and so continues from batch 10:
% it belongs to the run begun under the cover « Dualité de Cartier des
% groupes affines et groupes formels » (p. 170, batch 9), which batch 10
% carries to p. 200 (modules over W with F and V, A_F); the categories
% C_F, C'_F, C_1 and the objects A/FA, A/VA here are in that run's notation.
% Said in a \note at the head of p. 201. Pages
% 203-204 form one continuous text (203 ends « l'obstruction à l'obtenir à
% partir d'une », 204 begins « classe de faisceaux inversibles »); pages
% 205-207 are a second run, 206 continuing on 207 (« et considérons » /
% « lambda = Theta<z, z'> »), 207 breaking off at « on obtient ».
%
% What the batch is about. Page 201 (the tail of a run on formal groups /
% Dieudonné-type objects, reading largely doubtful): m/m^2 dual to the
% primitive elements P(U), finite dimension, finite socle, hence an
% « artinian » object; M/FM of finite dimension, finite length; a diagram of
% full subcategories C_0 -> C_F, C'_F -> C_1 -> \mathcal{C}, braced A_F and
% A_S with a circled « S = A - F », the isomorphism
% C_1/C_0 = C_F/C_0 x C'_F/C_0, the categories generated in C/C_0 by
% (A/FA)_{C_0} and by A/VA, and « limite projective / inductive d'objets
% artiniens ». Pages 203-204, « Accouplements de Tate »: A an abelian scheme
% over S, P a principal homogeneous bundle under A; Pic^*_{A/S} containing
% Pic^0_{A/S} = A', A acting trivially on A' and on NS_{A/S}; for P, the
% neutral component of Pic^*_{P/S} is canonically A', and the extension of
% NS_{A/S} by A' gives a class in Ext^1_S(NS_{A/S}, A'), hence a map from
% NS_{A/S} to H^1(S, A'); the obstruction to lifting a section lies in the
% Brauer group H^2(S, G_m); whence Tate's pairing
% H^1(S,A) x H^0(S,A') -> H^2(S,G_m), its refinement for S'/S faithfully
% flat, and for S' finite over S a hoped-for pairing
% \hat H^i(S'/S,A) x \hat H^{1-i}(S'/S,A') -> H^2(S'/S,G_m). Page 205,
% « Symboles locaux »: the tame symbol (f,g)_x on a projective non-singular
% curve over k algebraically closed, its properties (i)-(v) (bilinearity,
% (f,g)_x = f(x)^{v_x(g)}, triviality on strict units, the product formula,
% antisymmetry), the pairing on idèles I(X) x I(X) -> K^*, and the pairings
% K^* x C(X) -> k^* and K^*/k^* x C_0(X) -> k^* (a middle one struck).
% Pages 206-207, « la généralisation de Tate » (reading doubtful): X
% non-singular, A its Albanese variety, \hat A the Picard variety, N the
% Albanese kernel in Z_0(X), the group I of « répartitions », C_0 -> \hat A,
% and the construction of a pairing N x C_0 -> k^* by
% <z, \bar z> = Theta<z, z'> / prod_i f_{x_i}(x_i)^{a_i}, with Theta a divisor
% on X x \hat A (the X struck, see below), breaking off as he starts to show
% that Theta<z, z'> depends only on \dot D and z.
%
% Notation fixed once for the batch: \underline{\mathrm{Pic}} for his
% underlined Pic (as in folder 59); \mathrm{N.S.} kept with his dots;
% \mathbb{G}_m; \underline{\mathrm{Ext}}^1_S as he writes « Ext » with
% capital; his curly figure-3 for zero-cycles as \mathfrak{z}, and the group
% of zero-cycles as \mathfrak{Z}_0 (hand.md §4); the idèle/répartition group
% he writes with a script I as \mathrm{I}; vectors \vec f, \vec g as he
% arrows them; the local symbol (f,g)_x and the brackets <...> for
% « value on a cycle » set as \langle ... \rangle.
%
% Readings to check first: the whole of page 201 (fast, prose mostly lost);
% the sign after « de » on the cover; « neutre » inserted on p. 203; the word
% underlined after « Pic^*_{P/S} » on p. 203; the subscript « loc » on
% p. 204; the right-hand side « K^* » in (i) on p. 205 and the notes
% under K^*/k^*; « Tate » in the first line of p. 206; the phrase after
% « Sφ(z) » on p. 207.
%
% Apparatus count: 72 \ill{}, 39 \uncertain{}.
%
% Pass: Opus 5.5 (claude-opus-5-5), 2026-09-26 — first pass, unchecked against the pages by a human.

\documentclass[11pt,a4paper]{article}
% Le préambule commun à toutes les transcriptions du fonds Grothendieck.
%
% Il tient en une idée : **l'incertitude se compose, elle ne se résout pas.**
% Une transcription qui lisse un mot douteux a détruit la seule chose pour
% laquelle on la faisait — un lecteur doit pouvoir savoir, à chaque endroit,
% ce qui était sur le feuillet et ce qui a été ajouté. Les six macros qui
% suivent sont donc l'appareil critique, et non de la décoration.
%
% Les mêmes commandes sont reconnues par `scripts/render.mjs`, qui produit la
% vue de lecture du site. Une macro ajoutée ici doit l'être aussi là-bas,
% faute de quoi le rendu échouera bruyamment — ce qui est voulu : un rendu
% silencieusement faux est pire qu'un rendu absent.

\ProvidesPackage{grothendieck}[2026/08/08 Transcriptions du fonds Grothendieck]

\usepackage{amsmath,amssymb,amsthm}
\usepackage{mathrsfs}
\usepackage{stmaryrd}
% Les diagrammes commutatifs sont partout dans ce fonds : c'est la langue
% même de la géométrie algébrique de Grothendieck, et une transcription qui
% les rendrait en prose ne transcrirait rien.
\usepackage{tikz-cd}
% Français par défaut — c'est la langue des feuillets — mais l'anglais est
% chargé aussi : la traduction et le résumé sont des documents anglais, et la
% typographie française (espaces avant les deux-points) y serait une faute.
% Ces documents basculent par \selectlanguage{english} en tête de corps.
\usepackage[english,french]{babel}
\usepackage{fontspec}
\usepackage{geometry}
\geometry{a4paper,margin=25mm,marginparwidth=28mm}
\usepackage{marginnote}
\usepackage{xcolor}
% ulem, not soul. `soul`'s \st and \ul are built on a character-by-character
% scan that XeLaTeX's font machinery breaks: the document fails with an
% undefined control sequence rather than degrading. ulem does the same two jobs
% and is Unicode-engine safe. `normalem` stops it hijacking \emph, which the
% transcriptions use for Grothendieck's own emphasis.
\usepackage[normalem]{ulem}
\usepackage{hyperref}
\hypersetup{colorlinks=true,linkcolor=black,urlcolor=black,pdfborder={0 0 0}}

% --- Métadonnées du lot -----------------------------------------------------
% Reprises telles quelles par le script de rendu, qui les affiche en tête de
% la vue de lecture. Elles fixent ce que le fichier prétend couvrir : sans
% elles, une transcription flotte sans point d'attache dans un fonds de seize
% mille feuillets.

\newcommand{\folder}[1]{\gdef\grothFolder{#1}}
\newcommand{\batch}[1]{\gdef\grothBatch{#1}}
\newcommand{\leaves}[2]{\gdef\grothFirst{#1}\gdef\grothLast{#2}}
\newcommand{\foldertitle}[1]{\gdef\grothFoldertitle{#1}}
\gdef\grothFolder{?}\gdef\grothBatch{?}\gdef\grothFirst{?}\gdef\grothLast{?}\gdef\grothFoldertitle{}

% --- Le résumé --------------------------------------------------------------
%
% Il ouvre la lecture modernisée, sous le titre « Résumé », et il est composé
% en petit. Ce n'est pas un ornement : c'est le seul passage du document qui
% ne soit pas des mathématiques, et un lecteur doit pouvoir le reconnaître
% avant de l'avoir lu — puis le sauter s'il connaît déjà le sujet, ou s'y
% arrêter s'il ne le connaît pas. Le filet qui le suit marque la reprise du
% corps.

\newenvironment{resume}
  {\small\setlength{\parskip}{0.45em}}
  {\par\medskip\hrule\medskip}

% --- Mots-clés ---------------------------------------------------------------
%
% En anglais, en clôture du résumé de la lecture modernisée : le vocabulaire
% moderne sous lequel chercher ce que les feuillets construisent. Le manifeste
% du site (`npm run manifest`) les extrait pour étiqueter la cote entière —
% cette ligne est la source unique des tags, il n'y a pas de fichier à côté.
\newcommand{\keywords}[1]{\par\smallskip\noindent
  {\footnotesize\textsc{Keywords} --- \textit{#1}}\par}

% --- L'appareil critique ----------------------------------------------------

% Le numéro de feuillet, en marge. C'est l'ancre qui permet de régler un
% désaccord sur une formule en regardant *un* feuillet plutôt qu'une chemise
% de six cents. Le site s'en sert aussi pour tourner les pages du fac-similé
% quand on fait défiler la transcription.
\newcommand{\leaf}[1]{%
  \marginnote{\footnotesize\color{gray!70}\textsf{#1}}%
  \label{leaf:#1}\ignorespaces}

% Illisible : on ne devine pas. Un mot inventé qui se lit comme les autres est
% le pire résultat possible.
\newcommand{\ill}{\textcolor{red!60!black}{[\,\ldots\,]}}

% Lecture douteuse : le mot est donné, et signalé comme incertain.
\newcommand{\uncertain}[1]{\uline{#1}}

% Ajout de l'éditeur — une lettre manquante, un mot sous-entendu.
\newcommand{\add}[1]{\textcolor{blue!50!black}{[#1]}}

% Biffé par Grothendieck lui-même. Ce qu'il a rayé fait partie du document :
% une rature dit souvent où la pensée a changé de direction.
\newcommand{\struck}[1]{\sout{#1}}

% Note du transcripteur, sur l'état matériel du feuillet.
\newcommand{\note}[1]{\par\smallskip{\small\color{gray!60!black}\textit{[#1]}}\par\smallskip}

% Note marginale de Grothendieck, distincte du corps du texte.
\newcommand{\marginal}[1]{\par\smallskip\noindent\hspace{1em}%
  {\small\color{gray!60!black}#1}\par\smallskip}

% --- Titre ------------------------------------------------------------------

\AtBeginDocument{%
  \begin{center}
    {\large\bfseries Fonds Alexandre Grothendieck --- cote n\textsuperscript{o}~\grothFolder}\\[2pt]
    {\itshape\grothFoldertitle}\\[4pt]
    {\small feuillets \grothFirst--\grothLast\ (lot \grothBatch)}\\[2pt]
    {\footnotesize\color{gray!60!black}Transcription établie d'après le fac-similé numérisé
     par l'Université de Montpellier.\\
     Les lectures incertaines sont soulignées, les illisibles notées [\,\ldots\,],
     les ajouts entre crochets.}
  \end{center}
  \vspace{1em}}

\endinput


\folder{43}
\batch{11}
\pages{201}{207}
\dating{[à partir de 1961-vers 1968]}
\watermark{Édition de démonstration}
\foldertitle{Dualité des groupes. Jacobiennes généralisées etc : notes et copies de notes manuscrites (s.d.)}

\begin{document}

\page{201}

\note{la page commence au milieu d'une phrase : elle continue la suite
ouverte en page 170 sous la couverture « Dualité de Cartier des groupes
affines et groupes formels » et poursuivie jusqu'à la page 200 (modules sur
l'anneau de Dieudonné $A$, $F$ et $V$) ; aucune pagination de sa main.
Écriture rapide à l'encre bleue ; la prose de liaison se lit mal, les
formules et les mots soulignés mieux}

$\mathfrak{m}/\mathfrak{m}^2$ dual de $P(\mathfrak{U})$
\uncertain{primitifs}, \struck{\ill{}} \ill{} \ill{} \ill{} \ill{}
\ill{} \ill{} \ill{} par le \uncertain{noyau} de $F$.
\note{la lettre dans $P(\,\cdot\,)$ est un U majuscule arrondi, rendu
$\mathfrak{U}$ ; la lecture est douteuse}

S'il est de dim finie, le \emph{\uncertain{socle}} est fini, donc \ill{}
\ill{} \ill{} objet \emph{artinien} [N.B. l'\ill{}
\add{\uncertain{(indécomp.)}} de la catégorie des \uncertain{bigèbres}
\uncertain{de Dieudonné} \ill{} \uncertain{artiniennes} \ldots], donc on
a \ill{} \ill{} \struck{\ill{}} de \uncertain{Chevalley} \ill{} \ill{}
\ill{} de type fini sur $A$. De plus, la \uncertain{condition} \ill{}
\ill{} $M/FM$ est de dim finie sur \ill{} \ill{} de \emph{longueur
finie}.
\note{« (indécomp.) » est écrit au-dessus de la ligne, une accolade le
ramène avant « de la catégorie »}

\begin{tikzcd}[column sep=small, row sep=small]
  & \mathcal{C} & \\
  & C_1 \arrow[u] & \\
  C_F \arrow[ur] & & C'_F \arrow[ul] \\
  & C_0 \arrow[ul] \arrow[ur] \arrow[d, no head] & \\
  & 0 &
\end{tikzcd}

\note{schéma redessiné. Une grande accolade à gauche, marquée $A_F$, embrasse
$\mathcal{C}$, $C_1$ et $C_F$ ; une autre à droite, marquée $A_S$, embrasse
$\mathcal{C}$, $C_1$ et $C'_F$ ; à côté, entouré : « $S = A - F$ ».
$A_F$ est à rapprocher du $A_{\mathbf{F}}$ de la page 200. Sous
$C_0$ un trait vertical sans flèche descend vers un petit $0$. Le
$\mathcal{C}$ du haut est cursif, les autres $C$ droits}

\marginal{cat. \uncertain{formée} \struck{\ill{}} engendrée dans
$C/C_0$ par l'objet $(A/FA)_{C_0}$ \ill{}}\note{à gauche du schéma,
reliée par un trait à la flèche $C_0 \to C_F$ ; un $\omega$ isolé au-dessous}

\ill{} \ill{} (cat. \uncertain{pleine} \uncertain{engendr.})
\uncertain{engendrée} dans $C_1/C_0$ par \ill{} \ill{} \uncertain{objets}
\ill{} \ill{} de $C_1/C_0$, \ill{} $F$ est \struck{\uncertain{bijectif}}
\uncertain{surjectif} dans l'\uncertain{objet} \ill{} $C_0$.
\note{ce texte est relié par une accolade à la flèche $C_0 \to C'_F$}

\[
C_1/C_0 \simeq C_F/C_0 \times C'_F/C_0
\]

[objet \struck{\ill{}} engendré par $A/FA$

[objet \struck{\ill{}} engendré par $A/VA$

\emph{\ill{} objet \struck{\ill{}} \uncertain{simple}}
\note{à droite, un sigle biffé suivi de « $F_2$ » (lecture douteuse)}

\emph{\uncertain{limite} \uncertain{projective}} \emph{d'objets
\uncertain{artiniens}}

\emph{\uncertain{limite} \uncertain{inductive} d'objets
\uncertain{artiniens}}
\note{ces deux lignes, soulignées, sont écrites en bas à droite, de part et
d'autre de l'en-tête imprimé du papier ; un petit arc de cercle coupé par un
trait est tracé au-dessus}

\section*{Accouplements de \ill{} Tate -- Lang}

\note{de sa main, en haut à droite d'un feuillet orange de garde (page 202),
qui ne porte rien d'autre et ne reçoit pas de numéro de page. Le signe entre
« de » et « Tate » ressemble à un F barré}

\page{203}

\marginal{Accouplements de Tate}\note{au feutre orange, encadré, en haut à
droite}

$A$ schéma abélien sur $S$

$P$ fibré principal homogène sous $A$ \quad (schéma projectif)

$\underline{\mathrm{Pic}}^{*}_{A/S}$ schéma en groupes, contenant le schéma
abélien $\underline{\mathrm{Pic}}^{0}_{A/S} = A'$ schéma \emph{dual} de $A$.

$A$ opère sur lui-même par \uncertain{translations}, donc opère sur
$\underline{\mathrm{Pic}}^{*}_{A/S}$ et par suite sur $A'$, et opère
\emph{trivialement} sur $A' = \underline{\mathrm{Pic}}^{0}_{A/S}$ et
$\underline{\mathrm{Pic}}^{*}_{A/S}/\underline{\mathrm{Pic}}^{0}_{A/S} =
\mathrm{N.S.}_{A/S}$
\[
\Bigl[\ \text{donc}\quad A \times
\underbrace{\underline{\mathrm{Pic}}^{*}_{A/S}/
\underline{\mathrm{Pic}}^{0}_{A/S}}_{\text{Néron--Severi}}
\longrightarrow \underline{\mathrm{Pic}}^{0}_{A/S} = A'\ \Bigr]
\]
\note{le « $= A'$ » est écrit sous $\underline{\mathrm{Pic}}^{0}_{A/S}$,
avec un signe égal vertical}

\marginal{\ill{} $x + \lambda . \bar{x}$, \quad
$x \in \underline{\mathrm{Pic}}^{*}_{A/S}(T)$, $\lambda \in A(T)$,
$\bar{x} \in \mathrm{N.S.}(T)$}\note{en marge gauche ; les trois $\in$ sont
écrits verticalement sous $x$, $\lambda$, $\bar{x}$}

Considérons $\underline{\mathrm{Pic}}^{*}_{P/S}$, il a une « \emph{composante}
\emph{connexe} \add{\uncertain{neutre}} » $\underline{\mathrm{Pic}}^{0}_{P/S}$,
\struck{\ill{}} qui est \emph{can. isom.} à $\underline{\mathrm{Pic}}^{0}_{A/S}
= A'$ [car $A$, opérant trivialement] --- car
$\underline{\mathrm{Pic}}^{*}_{P/S}$ \ill{} \emph{\ill{}} à $P$ et à
$\underline{\mathrm{Pic}}^{*}_{A/S}$ sur lequel $A$ opère de la façon
naturelle. [il est \uncertain{extension} de $\mathrm{N.S.}_{A/S}$ par
$\underline{\mathrm{Pic}}^{0}_{A/S} = A'$, donc il lui correspond un élément
de $\underline{\mathrm{Ext}}^{1}_{S}(\mathrm{N.S.}_{A/S}, A')$, donc un
\ill{} \struck{\ill{}} \add{\ill{}} de $\mathrm{N.S.}_{A/S}$ \ill{} dans
$H^1(S, A')$, \struck{\ill{}} \ill{} \ill{} \ldots]
\note{« neutre » est écrit au-dessus de la ligne, entre les guillemets ; en
marge gauche, en face de la ligne de $\underline{\mathrm{Ext}}^1$, un sigle
biffé}

Si alors on a une section de $A'$, l'obstruction à l'obtenir à partir d'une

\page{204}

classe de faisceaux inversibles \struck{\uncertain{locaux}} sur $S$ est dans
le groupe de Brauer $H^2(S, \mathbb{G}_m)$ [se rappeler que
\[
H^{*}(S, \mathbb{G}_m) \underset{p}{\Longleftarrow}
H^p(S, \underline{H}^q_{\mathrm{loc}}(\mathbb{G}_m))
\]
d'où
\[
H^n(S, \mathbb{G}_m) \longrightarrow
H^0(S, \underline{H}^n_{\mathrm{loc}}(\mathbb{G}_m))
\quad \ldots\ ].
\]
\note{la flèche de la suite spectrale est double, avec $p$ en dessous.
L'indice « loc » des deux $\underline{H}$ est une lecture douteuse}

On obtient ainsi \emph{l'accouplement de Tate}
\[
\boxed{\;H^1(S, A) \times H^0(S, A') \longrightarrow H^2(S, \mathbb{G}_m)\;}
\]

Lequel d'ailleurs se précise (chaque fois que $S' \to S$ est fid. plat) en
des accouplements
\[
H^1(S'/S, A) \times H^0(S, A') \longrightarrow H^2(S'/S, \mathbb{G}_m)
\]

Il semble, quand $S'$ est \emph{fini} sur $S$, et sachant que les
$\hat{H}^i(S'/S, A)$ sont définis, qu'on puisse définir aussi des
accouplements
\[
\hat{H}^i(S'/S, A) \times \hat{H}^{1-i}(S'/S, A') \longrightarrow
H^2(S'/S, \mathbb{G}_m)
\]
qu'il \uncertain{faudrait} comprendre \ldots

\page{205}

\marginal{[Symboles locaux \ldots]}\note{encadré en haut à droite}

$X$ courbe projective non singulière connexe déf. sur $k$ alg. clos,
$K = k(X)$.

Si $f, g \in k(X)^{*}$, $x \in X$, on définit
\[
\boxed{\;(f,g)_x = \frac{f^{v_x(g)}}{g^{v_x(f)}}(x)\,
(-1)^{v_x(f)\,v_x(g)}\;}
\]

On a
\begin{itemize}
  \item[(i)] On a une application bilinéaire $K^{*} \times K^{*} \to
  K^{*}$\note{le dernier $K$ se lit majuscule, comme les deux premiers}
  \item[(ii)] $(f,g)_x = f(x)^{v_x(g)} = f\langle (g)_x \rangle$ si
  $f : X \to \mathbb{G}_m$ défini en $x$
  \item[(iii)] $(f,g)_x = 1$ si $g$ est une unité stricte, i.e.
  $g(x) = 1$
  \item[(iv)] $\displaystyle\prod_{x \in X} (f,g)_x = 1$
  \item[(v)] $(f,g)_x\,(g,f)_x = 1$
\end{itemize}
\note{une accolade réunit (i) à (v)}

\marginal{si \ill{}, \ill{} \ill{} $f : X \to \mathbb{G}_m$ défini
\ill{} une \emph{\ill{}} \ill{} \ill{} $g \in K^{*}$}\note{en marge gauche,
en face de (ii) et (iii)}

On en déduit un accouplement
\[
\mathrm{I}(X) \times \mathrm{I}(X) \longrightarrow K^{*}
\qquad (\mathrm{I}(X),\ \text{idèles})\ \text{par}
\]
\[
(\vec{f}, \vec{g}) = \prod_x (f_x, g_x)_x
\]

On a $(\vec{f}, \vec{g}) = 1$
\begin{itemize}
  \item si $\vec{f}$ est principal et $\vec{g}$ principal \struck{\ill{}}
  \item si $\vec{f}$ ou $\vec{g}$ est une unité stricte
\end{itemize}
\note{les deux conditions sont réunies par une accolade}

En \uncertain{particulier}, $K^{*}$ et $UK^{*}$ sont orthogonaux, d'où
des accouplements
\[
\boxed{\;K^{*} \times \mathcal{C}(X) \longrightarrow k^{*}\;}
\]
où $\mathcal{C}(X)$ = classes d'\emph{idèles} $= \mathrm{I}/UK^{*}$,
qui induit un accouplement
\[
\{f, \vec{g}\} = f\langle D \rangle
\prod_{x \in \ldots} g_x(x)^{-v_x(f)}
\]
\note{l'ensemble d'indices du produit est illisible, rendu $\ldots$}
si \ill{} $\vec{g} \in \mathcal{C}(X)$ \ill{} \uncertain{principal} en $x$,
$D$ \ill{} la classe du diviseur de $\vec{f}$.
\note{« classes d'idèles » est écrit sous $\mathcal{C}(X)$ ; dans
« $\mathrm{I}/UK^{*}$ » une lettre avant $K^{*}$ est récrite. La formule de
$\{f, \vec{g}\}$ et la phrase qui la suit sont écrites à cheval sur deux
cadres, la phrase dans la marge droite}

\struck{$\mathcal{C}(X) \times UK^{*} \longrightarrow k^{*}$}
\note{encadré, barré en zigzag ; au-dessous, biffé : « classes d'idèles »}

\[
\boxed{\;K^{*}/k^{*} \times C_0(X) \longrightarrow k^{*}\;}
\]
\note{sous $K^{*}/k^{*}$ : « diviseurs \ill{} \ill{} » ; sous $C_0(X)$ :
« classes d'idèles de degré $0$ »}

\page{206}

Voici la généralisation de \uncertain{Tate}

$X$ variété non singulière déf. sur $k$ alg. clos, $K = k(X)$

$A = A(X)$ var. d'Albanese, \add{hom. $\varphi : X \to A(X)$}
$\hat{A}$ variété de Picard
\note{« hom. $\varphi : X \to A(X)$ » est écrit au-dessus de la ligne, avec
une flèche qui le ramène entre les deux variétés ; le « $(X)$ » de
$A(X)$ y est surchargé, et un signe biffé précède $\hat{A}$}

Soient $N$ le noyau d'Albanese : $N = \mathrm{Ker}(\mathfrak{Z}_0(X)
\to A(X))$

\struck{\ill{}} $\mathrm{I}$ le groupe des classes de
\uncertain{répartitions}
\[
\vec{f} = (f_x)_{x \in X}
\]
telles que il existe un diviseur $D$ sur $X$, avec
\[
(f_x)_x = [D]_x \quad \text{pour tout } x.
\]
On a une application
\[
C = \mathrm{I}/K^{*} \longrightarrow \mathrm{Pic}(X)
\quad (\text{classes de diviseurs}),
\]
soit $C_0$ \struck{\ill{}} l'image inverse de $\hat{A}$ ; donc on a un
\uncertain{homomorphisme} surjectif
\[
C_0 \longrightarrow \hat{A}
\]
[$\mathrm{I}_0$ est l'image inverse de $C_0$ dans $\mathrm{I}$].
\note{dans $\mathrm{I}/K^{*}$, un $k^{*}$ est écrit puis biffé avant
$K^{*}$. Le but de la flèche est surchargé ; on lit un P, rendu
$\mathrm{Pic}(X)$ d'après la parenthèse}

On veut définir un accouplement
\[
N \times C_0 \longrightarrow k^{*}
\]
et pour cela, on définit un accouplement
\[
N \times \mathrm{I}_0 \longrightarrow k^{*}
\]
\uncertain{nul sur} $N \times K^{*}$. Soit \add{$\bar{\mathfrak{z}} \in C_0$,}
\[
\mathfrak{z} = \sum a_i x_i \in N, \qquad
\vec{f} = (f_x)_{x \in X} \in \mathrm{I}_0
\]
un représentant de $\bar{\mathfrak{z}}$ tel que $\mathrm{div}(\vec{f})$ ne
\uncertain{rencontre} pas $\mathfrak{z}$,
$D$ le diviseur défini par $\vec{f}$, \struck{\ill{}} $\dot{D}$ sa classe
dans $\hat{A}$, et $\mathfrak{z}' \in \mathfrak{Z}_0(\hat{A})$,
\struck{\ill{}} $\deg \mathfrak{z}' = 0$, tel que $\dot{D} =
S(\mathfrak{z}')$. Soit \uncertain{enfin} $\Theta$ un diviseur
\uncertain{canonique} sur \struck{$X$} $\times \hat{A}$ \uncertain{ne
rencontrant} pas $\varphi(\mathfrak{z}) \times{}$ \struck{\ill{}}
$\mathfrak{z}'$, et considérons
\note{« $\bar{\mathfrak{z}} \in C_0$, » est écrit au-dessus de la ligne ;
sous $x_i$ un mot biffé. Le $X$ de $X \times \hat{A}$ est barré d'un trait
en croix, sans rien au-dessus ; même chose en page 207}

\page{207}

\[
\lambda = \Theta\langle \mathfrak{z}, \mathfrak{z}' \rangle
\]
\note{un signe biffé devant le premier $\mathfrak{z}$, un autre après le
crochet}
qui est défini en notant que $\Theta\langle \mathfrak{z} \rangle$ est un
\struck{\ill{}} diviseur sur $\hat{A}$ qui est linéairement équivalent à
$0$, car sa classe est définie par $S\varphi(\mathfrak{z})$ \ill{}, qui
est nulle ; donc on peut \struck{prendre} trouver une fonction $g$ sur
$\hat{A}$ telle que
\[
(g) = \Theta\langle \mathfrak{z} \rangle,
\]
et on pose
\[
\lambda = g\langle \mathfrak{z}' \rangle \qquad
(\text{indépendant de } g \text{ car } \deg \mathfrak{z}' = 0)
\]

D'autre part, on considère pour tout $x \in{}$ \struck{$X$} \ill{}
$\mathfrak{z}$ les \struck{\ill{} $f_x(x)$} $f_{x_i}(x_i)^{a_i}$, et on
pose
\[
\mu = \prod_i f_{x_i}(x_i)^{a_i}
\]
on pose alors
\[
\langle \mathfrak{z}, \bar{\mathfrak{z}} \rangle = \lambda/\mu =
\Theta\langle \mathfrak{z}, \mathfrak{z}' \rangle \Big/
\prod_i f_{x_i}(x_i)^{a_i}
\]

Montrons que $\Theta\langle \mathfrak{z}, \mathfrak{z}' \rangle$ ne dépend
que de $\dot{D}$ et $\mathfrak{z}$ : si on change $\Theta$ en
$\Theta + (h)$, ($h$ fonction sur \struck{$X$} $\times \hat{A}$)
\struck{\ill{}} on obtient
\note{la page s'arrête ici, au premier tiers}

\end{document}
